\chapter{Sửa cỗ máy khi không có sơ đồ}
% full version: chapters/17_debugging.tex

\pencilfig{17}{\pencilart{17}}{Bo mạch không có sơ đồ: đầu dò chỉ nghe được vài dây dẫn trong hàng tỷ.}

Hãy tưởng tượng bạn được giao một bo mạch đang chạy mà không có sơ đồ, và được yêu cầu sửa nó. Kỹ sư có bốn công cụ: đầu dò để nghe tín hiệu; máy phát để đưa tín hiệu của mình vào; khả năng ngắt một linh kiện để xem cái gì hỏng; và quyền truy cập các thanh ghi điều khiển. Người ta nghiên cứu bộ não bằng đúng những công cụ đó.

\section*{Đầu dò cho một nghìn nơ-ron}

Một điện cực cắm vào não nghe được tiếng tách của vài nơ-ron ở gần. Tưởng như so với tám mươi sáu tỷ, một nghìn điện cực chẳng thấm vào đâu. Nhưng chừng ấy đủ để điều khiển con trỏ trên màn hình. Vì sao? Các nơ-ron lân cận nhiễu giống nhau, và hoạt động điều khiển bàn tay thực ra được mô tả chỉ bằng một, hai chục biến chung. Không cần nghe tất cả, chỉ cần nghe thấy vài ``giọng nói'' này.

Tín hiệu thô từ một nghìn điện cực là hàng trăm triệu bit mỗi giây: kênh vô tuyến từ dưới hộp sọ không kham nổi. Còn bản thân các tiếng tách mang ít hơn khoảng một nghìn lần. Vì thế nên tách các tiếng tách ngay trên thiết bị cấy và chỉ gửi chúng ra ngoài --- cũng là mẹo nén như trong mắt.

\section*{Neuralink làm gì}

Công ty Neuralink nhận phần kỹ thuật của bài toán: những sợi mảnh, dẻo thay cho kim cứng, một robot cấy chúng, và một thiết bị cấy không dây kín khít. Công trình họ công bố mô tả một hệ thống nối với máy tính bằng cáp, còn việc nén và truyền không dây được nêu là kế hoạch. Dữ liệu về người bị liệt tay điều khiển con trỏ hiện chủ yếu được biết qua báo cáo của chính công ty; theo đó, thiết bị cấy có khoảng một nghìn kênh trên vài chục sợi, và một phần các sợi đã bị xê dịch sau khi cấy. Bản thân ý tưởng không mới: các nhóm học thuật dùng mảng điện cực cứng khoảng một trăm điện cực đã giúp con người điều khiển con trỏ, ``viết'' chữ bằng ý nghĩ và biến những lần cố nói thành văn bản trên màn hình với tốc độ khoảng sáu mươi từ mỗi phút.

\section*{Đọc và ghi}

Bộ giải mã đọc tần số tiếng tách của nhiều nơ-ron và thu được chưa đến mười bit mỗi giây --- một phần nhỏ nhoi so với những gì bản thân các nơ-ron mang. Bù lại, cả bộ não lẫn bộ giải mã học lẫn nhau: con người thích nghi với bộ giải mã, bộ giải mã --- với con người. Hai học trò tự điều chỉnh theo nhau là một bài toán tinh tế: ngay cả khi mỗi bên riêng lẻ đều ổn định, cùng nhau chúng vẫn có thể dao động.

Ghi vào não khó hơn đọc. Dòng điện qua điện cực kích thích mọi thứ xung quanh chứ không phải một nơ-ron cần thiết. Có thể ``vẽ'' những điểm sáng, như các điểm ảnh phát sáng, và người đó sẽ phân biệt được các chữ cái đơn giản. Nhưng ghi mã nén của mắt sao cho bộ não nhìn thấy hình ảnh thì đến nay chưa ai làm được.

\section*{Điều đầu dò không thấy}

Điện cực nghe được mạng điện thoại truyền dữ liệu. Nhưng các đài phát thanh --- dopamine, serotonin, noradrenaline, acetylcholine --- là nồng độ các chất, đầu dò không nghe được chúng. Nó cũng không thấy độ mạnh của các liên kết, nơi lưu toàn bộ trí nhớ. Và không thấy cơn đau theo cách con người cảm nhận nó. Một trình gỡ lỗi thấy dữ liệu nhưng không thấy các thanh ghi sẽ luôn ngạc nhiên vì sao cùng một đầu vào lại cho những câu trả lời khác nhau.

\fullversion{17}
